\section{Related Work}

\label{sec:related_work}

Lexical retrieval methods rank documents by term frequency and inverse document frequency with length normalization. BM25, the probabilistic lexical ranker used as our primary baseline, has remained competitive with learned retrieval methods on many benchmarks because of its simplicity and robustness. The agent memory setting differs from classical IR in one structural way: candidate documents are written by the agent itself during task execution, giving them timestamps and access frequencies that can serve as auxiliary ranking signals alongside lexical overlap. Lost in the Middle demonstrates that relevant information position affects multi-document question answering, motivating attention to which items reach the prompt and how they are ordered within it \citep{ref_9dbf664b6954efbc}.

Agent memory architectures have been proposed at increasing levels of autonomy. MemGPT introduces an operating-system-inspired virtual context manager that moves information across memory tiers using explicit management calls, demonstrating document analysis and multi-session chat scenarios \citep{ref_959d6858d034c6e1}. The present study draws on the conceptual vocabulary of tiered storage but does not implement its autonomous virtual-memory controller. A-MEM organizes memories as structured notes linked by dynamically computed associations following a Zettelkasten-inspired design \citep{ref_1aacba116cbe3c5b}; our fixed lexical ranking lacks autonomous linking and evolution. Generative Agents synthesize natural-language experience records into reflections and dynamically retrieve them for planning behavior in interactive simulation \citep{ref_6e953134a3c04402}; we draw on the retrieval framing but implement neither reflection nor planning. MemoryBank supports retrieval and updates for sustained personalized interaction, modeling forgetting and reinforcement influenced by elapsed time and importance \citep{ref_713e886950406893}; our temporal-decay term is a simpler heuristic.

LongMemEval evaluates five core long-term memory abilities -- information extraction, multi-session reasoning, temporal reasoning, knowledge updates, and abstention -- and explicitly separates indexing, retrieval, and reading stages \citep{ref_ab2c226739f2f52a}. The controlled evaluation in this paper uses a fixed-hash balanced subset with disclosed deviations including a substituted judge and a segmented bank; it is not an official benchmark result. The specific gap our work addresses -- the interaction between query-independent temporal weights and lexical relevance under a fixed budget -- has not been analyzed formally or addressed with a provable structural fix in prior work on agent memory selectors.
